\documentclass[11pt]{article}
\usepackage[margin=24mm]{geometry}
\usepackage{amsmath,amssymb,amsthm,mathtools}
\newcommand{\Q}{\mathbb Q}
\newcommand{\Gr}{\operatorname{Gr}}
\newcommand{\W}{\operatorname{W}}
\newcommand{\HH}{\mathrm H}
\newtheorem{theorem}{Theorem}
\newtheorem{corollary}[theorem]{Corollary}
\newtheorem{lemma}[theorem]{Lemma}
\title{GST genuine causal topology folds on the original arbitrary X:
all-N emergent mixed Hodge weights from graph incidence without dimension elevation}
\date{21 October 2026}
\begin{document}\maketitle

\section{Corrected interpretation: dimension is born from a GST realization, not a coordinate shifted after birth}
The user's actual proposal starts from the
\emph{dimensionless GeneralSpace ontology},
whose points and associative paths precede topology,
Krull dimension, complex structure and cohomology.
It asks for a FOLD OF THE REALIZED TOPOLOGY ITSELF,
controlled by GST higher-causality, that forces a
specific equation deciding the presence or absence
of Hodge structure on ANY original X.
This is NOT the preceding experiment
\(\mathbf P_X(E)\), which changes an auxiliary fiber
dimension while split-preserving all original
Hodge/Chow defect sectors.
It is NOT a nonholomorphic self-map of a fixed
elliptic product, which rearranges classes but leaves
the topology unchanged.

The actual source-file interpretation is:
\texttt{GSTGeneralSpace.lean} has no dimension field;
\texttt{GSTGeneralSpaceSchemeDimensionObservation.lean}
extracts Krull dimension as an actual scheme observation;
\texttt{GSTGraphV2SelfExpandingCosmology.lean}
allows a new world born from arbitrary-arity parent
worlds with new internal law types;
\texttt{GSTGraphV2RecursiveStateCausality.lean}
has local and cross-world causal paths;
\texttt{GSTGraphV2UnboundedHigherCausality.lean}
has source/target boundary equations in every degree.
These abstract laws do NOT automatically provide
topological quotient schemes: below we construct
such schemes ACTUALLY, over every complex X.

\section{An N-fold ACTUAL intrinsic topology fold of any original smooth projective X}
Fix ANY geometrically connected smooth projective
complex variety X of dimension d>=1, and any N>=0.
Choose 2N distinct closed complex points
\[
 p_1,q_1,\ldots,p_N,q_N\in X(\mathbb C).
\]
This is possible for EVERY N when d>=1.
Construct the actual finite algebraic pinching quotient
\[
 \boxed{
 \nu_N:X\longrightarrow Y_N
        =X/(p_1\sim q_1,\ldots,p_N\sim q_N).
 }                                                        \tag{1}
\]
It is a projective complex algebraic scheme;
the normal source X is its normalization, and the
N identified target points are isolated singularities.
The map is finite, proper and birational.
Its underlying geometric dimension is UNCHANGED:
\[
 \boxed{\dim_{\mathbb C}Y_N=d.}                       \tag{2}
\]
This operation folds the topology of X itself;
it does NOT multiply X by another auxiliary variety.

The exact ordinary rational-constant-sheaf
normalization sequence is
\[
 0\longrightarrow \Q_{Y_N}
 \longrightarrow (\nu_N)_*\Q_X
 \xrightarrow{\mathrm{diff}}
 \bigoplus_{e=1}^N\Q_{r_e}
 \longrightarrow 0,                                  \tag{3}
\]
where at the e-th pinched point r_e the difference
map is \(a_e-b_e\) on its two normalization branches.
For a finite proper map with finite discrete fibers,
the higher direct images of \(\Q_X\) vanish,
so taking sheaf cohomology is legitimate.
X is connected: the H^0 map from Q to Q^N
in the long exact sequence is identically ZERO
(each global constant has equal values at p_e,q_e).

\begin{theorem}[UNIVERSAL intrinsic GST N-fold causal topology birth theorem]
For EVERY connected smooth projective complex
d-fold X with d>=1 and EVERY integer N>=0,
the genuine folded projective algebraic scheme Y_N
has the same Krull dimension d but its first
cohomology sits in the canonical short exact sequence
of rational mixed Hodge structures
\[
 \boxed{
 0\longrightarrow\Q(0)^N
 \longrightarrow H^1(Y_N,\Q)
 \xrightarrow{\nu_N^*}H^1(X,\Q)
 \longrightarrow0.
 }                                                        \tag{4}
\]
Therefore its mixed-Hodge weight-zero subspace and
weight-one graded piece satisfy
\[
 \boxed{
  W_0H^1(Y_N,\Q)\cong\Q(0)^N,\quad
  \Gr^W_1H^1(Y_N,\Q)\cong H^1(X,\Q),\quad
  b_1(Y_N)=b_1(X)+N.
 }                                                        \tag{5}
\]
For k>=2 the normalization induces an isomorphism
of rational mixed Hodge structures
\(H^k(Y_N,\Q)\cong H^k(X,\Q)\).
In particular, for EVERY N>0 the projective
Y_N does NOT have pure weight-one H^1,
even when the original X had entirely pure
smooth-projective Hodge structure.
The induced change is a consequence of TRUE
topology folding of the ORIGINAL X, with
NO change of its algebraic dimension.
\end{theorem}
\begin{proof}
Use the exact finite-normalization sheaf
sequence (3) and the fact that N distinct
singular points contribute a skyscraper
coefficient \(\Q^N\) in degree zero and no
higher sheaf cohomology.
Since X and Y_N are connected, the map on
global constant functions is an isomorphism
\(\Q\to\Q\), while the map to \(\Q^N\)
vanishes identically.
The long exact sequence yields
\(0\to\Q^N\to H^1(Y_N)\to H^1(X)\to0\)
and isomorphisms in degrees k>=2.
The constant sheaf normalization triangle
is a morphism in the category of
constructible mixed Hodge complexes of
proper complex algebraic varieties;
its H^0 skyscraper term has pure
weight 0, while H^1(X) is pure weight 1
because X is smooth projective.
Strictness of the canonical mixed Hodge
weight filtration yields exactly (4),(5).
Finite birational maps preserve Krull
dimension, establishing (2).
No Hodge conjecture hypothesis was
used in the proof.
\end{proof}

\begin{corollary}[Actual same-dimensional Hodge breaking on every projective X]
The N-fold topology fold fails to preserve
PURE Hodge structure if and only if N>0.
For the boundary case X=P^d (d>=1),
where H^1(X,Q)=0, it gives the EXACT,
strongest simple equation
\[
 \boxed{H^1(Y_N,\Q)\cong\Q(0)^N.}                 \tag{6}
\]
Thus even varieties with completely
algebraic even-degree cohomology acquire
new first cohomology from causal topology alone.
The new H^1 is of weight 0, and so it is
NOT the H^1 of a smooth projective variety.
The result does NOT contradict the
rational Hodge conjecture: the new Y_N
is singular and the relevant assertion is
a mixed-weight anomaly, not a nonalgebraic
rational (p,p) Hodge class on smooth X.
\end{corollary}

\section{Arbitrary finite GST CAUSAL GRAPH fold: Hodge weight zero exactly equals loop rank}
The previous construction is the bouquet graph
with one vertex (the original connected X)
and N independent causal loop identifications.
There is a more general ACTUAL algebraic folding
law across arbitrary smooth projective input worlds.
Let \(\Gamma=(V,E,s,t)\)
be a finite directed multigraph.
For each v choose a connected smooth projective
complex variety \(X_v\) of dimension >=1.
Choose disjoint actual closed points at the ends
of every oriented graph edge e, and identify the
corresponding two points. Denote the resulting
actual projective scheme by \(Y_\Gamma\)
and its normalization by \(\bigsqcup_v X_v\).
Let \(c\) be the number of connected components
of the underlying undirected graph.
Its rational cellular 1-coboundary is
\[
 \boxed{\delta_\Gamma:\Q^{V}\to\Q^E,\qquad
 (\delta_\Gamma f)(e)=f(t(e))-f(s(e)).}           \tag{7}
\]
The entire influence of the gluing in first
cohomology is governed by this ACTUAL
causal incidence matrix.

\begin{theorem}[GST graph topology-fold equation, ALL graphs and ALL projective input dimensions]
For every such actual graph gluing,
\[
 \boxed{
 0\longrightarrow\operatorname{coker}\delta_\Gamma
 \longrightarrow H^1(Y_\Gamma,\Q)
 \longrightarrow \bigoplus_{v\in V}H^1(X_v,\Q)
 \longrightarrow0
 }                                                        \tag{8}
\]
is an exact sequence of rational mixed Hodge structures.
Furthermore
\[
 \boxed{
 \Gr^W_0 H^1(Y_\Gamma,\Q)
   \cong H^1(\Gamma,\Q),\qquad
 \dim_\Q \Gr^W_0H^1(Y_\Gamma,\Q)
   =|E|-|V|+c.
 }                                                        \tag{9}
\]
Every connected graph has pure weight-one H^1
if and only if its incidence graph is a TREE;
if it has a loop, H^1 has a nonzero
weight-zero sector and cannot be pure.
The conclusion is independent of the algebraic
dimensions of all the \(X_v\) (provided >=1).
\end{theorem}
\begin{proof}
The normalization sheaf sequence is
\(0\to\Q_Y\to \nu_*\Q_{\bigsqcup X_v}
\to\oplus_e \Q_{r_e}\to0\).
Taking cohomology, the map
\(H^0(\bigsqcup_v X_v)=\Q^V
\to \oplus_eH^0(\{r_e\})=\Q^E\)
is precisely the signed edge incidence
\(\delta_\Gamma\), up to a harmless
overall sign convention.
The same long exact sequence gives (8);
the skyscraper term has weight zero and
each smooth-projective H^1(X_v) has
pure weight one.
The kernel of \(\delta_\Gamma\)
has dimension c (functions constant on
graph connected components), so
\(\operatorname{rank}\delta_\Gamma=|V|-c\).
Thus its cokernel has dimension
\(|E|-|V|+c\), which is precisely
the cellular first Betti number
of \(\Gamma\).
For a connected graph this vanishes
exactly when \(\Gamma\) has no graph
cycles, i.e. when it is a tree.
\end{proof}

\section{The higher-GST CAUSAL BOUNDARY law: why source/target coherences are algebraically decisive}
The actual
\texttt{GSTGraphV2UnboundedHigherCausality.lean}
has globular identities
\[
 s(sx)=s(tx),\qquad t(sx)=t(tx)
 \quad \text{for every \((n+2)\)-cell }x.
\]
Linearize its genuine cell sets over Q:
\(C_n^{GST}=\Q[H.\mathrm{Cell}(n)]\),
and define signed source-minus-target
boundary on each \((n+1)\)-cell:
\[
 \boxed{\partial_{n+1}[x]=[s(x)]-[t(x)].}          \tag{10}
\]
These are actual equations implied by the
EXISTING higher-causality structure.

\begin{theorem}[Globular GST higher-causality is an authentic chain complex]
For every index n>=0,
\[
 \boxed{\partial_{n+1}\partial_{n+2}=0.}          \tag{11}
\]
Thus all finite higher-causal dimensions can be
made into a rational chain complex.
For a finite causal 1-skeleton with b1 loops,
the higher causal 2-cells may impose linearly
independent relations; if their signed
boundaries span an r-dimensional subspace
of the kernel of the 1-boundary, then
\[
 \boxed{\dim_\Q H_1^{GST}=b_1-r.}                 \tag{12}
\]
This algebraic law is an EXACT way to
decide which topology-loop classes remain
after enforcing second-order causal relations.
Its identification with weight-zero
H^1 of an ALGEBRAIC variety requires
a further actual normalization/hypercover
realization; arbitrary globular 2-cells
do not automatically fill algebraic loops.
\end{theorem}
\begin{proof}
Expand
\[
 \begin{aligned}
 \partial_{n+1}\partial_{n+2}[x]
 &=[s(sx)]-[t(sx)]-[s(tx)]+[t(tx)]\\
 &=0
 \end{aligned}
\]
by the two literal globular boundary equations.
For a finite chain complex,
\(H_1=\ker\partial_1/\mathrm{im}\partial_2\)
and \(\mathrm{im}\partial_2\subset\ker\partial_1\).
Subtracting the rank of \(\partial_2\)
proves (12). The comparison with actual
mixed Hodge cohomology is warranted only
for geometrically certified cells;
equations (8),(9) give that certification
for point-pinch graph 1-skeleta.
\end{proof}

\section{Genuinely ALL-N higher-degree causal weight birth, beyond H^1}
Let C be the actual projective
nodal rational curve with normalization P^1
and one point-pair pinched. Its true rational
mixed Hodge cohomology is
\[
 H^0(C)=\Q(0),\quad
 H^1(C)=\Q(0),\quad
 H^2(C)=\Q(-1).
\]
For ANY original smooth projective complex
X of arbitrary dimension and ANY N>=1,
form the actual N-direction folded world
\[
 Y_{X,N}=X\times C^N
 \quad (\dim_{\mathbb C}Y_{X,N}=\dim X+N).
\]
Its entire rational mixed Hodge cohomology
is the genuine Kunneth tensor decomposition
\[
 \boxed{
 H^*(Y_{X,N},\Q)
   \cong H^*(X,\Q)\otimes
  (\Q(0)[0]\oplus\Q(0)[1]\oplus
                       \Q(-1)[2])^{\otimes N},
 }                                                        \tag{13}
\]
where brackets give cohomological-degree
positions, not Tate weights.
Let \(\gamma_i\in H^1(C,\Q)=\Q(0)\)
be the independently born causal-loop classes.
For ANY nonzero rational \((p,p)\) Hodge class
\(\alpha\in H^{2p}(X,\Q)\), the ACTUAL
class
\[
 \boxed{
 \beta_N=\pr_X^*\alpha\smile
      \gamma_1\smile\cdots\smile\gamma_N
 \in H^{2p+N}(Y_{X,N},\Q)
 }                                                        \tag{14}
\]
is NONZERO for every N, by the
Kunneth decomposition.
Its Hodge type as a mixed-Hodge class
is \((p,p)\), its WEIGHT is \(2p\),
while its COHOMOLOGICAL DEGREE is \(2p+N\):
\[
 \boxed{
 \operatorname{wt}(\beta_N)=
     \deg(\beta_N)-N.
 }                                                        \tag{15}
\]
This is an exact ALL-N law that makes
causally folded topological dimensions
emerge as a mismatch between
cohomological degree and Hodge weight.
It needs no rational Hodge conjecture,
only an actual rational Hodge class alpha
on the original X (the unit class
alpha=1 always qualifies for p=0).
It does NOT give a rational Hodge counterexample
on a smooth projective variety:
the folded Y is singular for N>=1,
so mixed weights are the correct invariants.

\section{The sharp logical test of the user's hypothetical full-Hodge premise}
Suppose hypothetically that the rational
Hodge conjecture holds for every
smooth projective complex X.
The intrinsic GST point-pinch theorem
still produces (4) for ALL X,d,N:
there is no contradiction because
\(Y_N\) is singular, and the new
weight-zero H^1 is not an unrepresented
rational \((p,p)\) Hodge class on smooth X.
A genuine contradiction would require
a fold output satisfying ALL THREE:
(i) the output Z is smooth projective,
(ii) the output carries a genuine rational
\((p,p)\) Betti class,
(iii) that class is proven NOT in the
canonical rational Chow image.
Neither unbounded causal coherence nor
the mixed-Hodge rank equation asserts (iii).

What has changed is significant and precise:
the new equations genuinely modify
topology and CREATE cohomology on the
original X of unchanged observed dimension.
The Graph V2 causal loop rank
\[
       |E|-|V|+c
\]
directly governs the birth of a nonpure
mixed Hodge sector. This is the first
dimension-free GST graph-fold mechanism
with an actual all-X, all-N
scheme-theoretic mixed-Hodge realization,
as opposed to a class-only transformation.
\end{document}
